\paragraph{Kebutuhan Nonfungsional}

Selain kebutuhan fungsional, sistem audit trail juga perlu memenuhi sejumlah kebutuhan nonfungsional yang berkaitan dengan kualitas layanan, khususnya mengingat sifat audit record sebagai bukti digital yang harus terjamin integritas, ketersediaan, dan keamanannya. Kebutuhan nonfungsional disusun mengacu pada karakteristik kualitas perangkat lunak ISO/IEC 25010 serta kontrol AU (\textit{Audit and Accountability}) pada NIST SP 800-53 Rev. 5 \cite{NIST80053}. Daftar kebutuhan nonfungsional disajikan pada Tabel~\ref{tab:kebutuhan-nonfungsional}.

\begin{xltabular}{\textwidth}{|>{\centering\arraybackslash}p{1.4cm}|>{\raggedright\arraybackslash}p{2.6cm}|X|}
    \caption{Kebutuhan Nonfungsional Sistem Audit Trail} \\
    \hline
    \textbf{ID} & \textbf{Kategori} & \textbf{Deskripsi} \\
    \hline
    \endfirsthead
    \hline
    \textbf{ID} & \textbf{Kategori} & \textbf{Deskripsi} \\
    \hline
    \endhead
    \hline
    \endfoot
    \hline
    \endlastfoot

    NF01 & Kinerja & Sistem harus mampu membangkitkan dan menyimpan audit record dengan tambahan latensi yang tidak mengganggu waktu respons proses utama DecMed. \\
    \hline

    NF02 & Kinerja & Sistem harus mampu menampilkan hasil pencarian audit record dalam waktu respons yang wajar meskipun volume audit record yang tersimpan besar. \\
    \hline

    NF03 & Skalabilitas & Sistem harus mampu menangani pertumbuhan volume audit record seiring meningkatnya jumlah transaksi pada jaringan IOTA tanpa penurunan kinerja yang signifikan. \\
    \hline

    NF04 & Ketersediaan & Layanan pembangkitan audit trail harus tersedia secara konsisten agar tidak ada aktivitas penting pada sistem DecMed yang gagal tercatat. \\
    \hline

    NF05 & Keandalan & Sistem harus menjamin audit record tidak hilang meskipun terjadi gangguan sementara pada komponen pengumpul (\textit{collector}), misalnya melalui mekanisme antrean atau \textit{buffer}. \\
    \hline

    NF06 & Integritas & Audit record harus dilindungi dari modifikasi yang tidak sah melalui mekanisme kriptografis (\textit{hash chaining}), dan pelanggaran integritas harus dapat terdeteksi. \\
    \hline

    NF07 & Keamanan & Akses terhadap audit trail harus dilindungi melalui autentikasi dan otorisasi berbasis peran, serta audit record yang bersifat sensitif harus terenkripsi baik saat disimpan (\textit{at rest}) maupun saat dikirimkan (\textit{in transit}). \\
    \hline

    NF08 & Daya Tahan (\textit{Durability}) & Audit record harus tersimpan secara tahan lama sesuai kebijakan retensi yang ditentukan, tanpa bergantung pada mekanisme \textit{pruning} penyimpanan eksternal seperti pada jaringan IOTA. \\
    \hline

    NF09 & Usabilitas & Antarmuka pencarian audit trail harus dapat digunakan oleh auditor tanpa memerlukan keahlian teknis mendalam terhadap struktur data \textit{distributed ledger}. \\
    \hline

    NF10 & Kepatuhan & Struktur dan pengelolaan audit record harus sesuai dengan pedoman kontrol AU pada NIST SP 800-53 Rev. 5 serta ketentuan perlindungan data pribadi yang berlaku, yaitu Undang-Undang Nomor 27 Tahun 2022 tentang Pelindungan Data Pribadi. \\
    \hline

    NF11 & Interoperabilitas & Audit record harus dapat diekspor dalam format standar (mis. JSON atau CSV) agar dapat diintegrasikan dengan sistem pemantauan keamanan lain. \\
    \hline

    NF12 & Pemeliharaan (\textit{Maintainability}) & Desain sistem audit trail harus modular sehingga penambahan jenis \textit{event} baru di masa mendatang tidak memerlukan perubahan besar pada arsitektur yang sudah ada. \\
    \hline

\end{xltabular}